% Pista ana-23j-q5 · Anàlisi
% Procedència: PAU Catalunya 2023 · Convocatòria de juny · Sèrie 1 · Qüestió 5
\documentclass[11pt,a4paper]{article}
\usepackage{pista}
\pistainfo{Catalunya 2023}{Convocatòria de juny}{Sèrie 1}

\begin{document}

\section*{\textcolor{blauPista}{Qüestió 5}}

\noindent La Núria té un jardí rectangular i vol fer-hi un tancat (rectangular o quadrat) de $8\,\text{m}^{2}$ per al seu gos. Ha pensat de posar el tancat tocant al mur del jardí, tal com es mostra a la figura de la dreta, per estalviar-se així un dels quatre costats. El preu de la tanca que vol fer servir és de $2{,}5\,\euro/\text{m}$.

\subsection*{\textit{a)} \normalfont Quines dimensions ha de tenir el tancat perquè el cost sigui mínim? Quin és aquest cost mínim? \hfill\textnormal{[1,75~punts]}}

\pista{És un problema d'optimització.
\begin{itemize}
\item \textit{Variables.} Anomena $x$ el costat paral·lel al mur i $y$ el perpendicular.
\item \textit{Condició.} L'àrea, $x \cdot y = 8$, et permet expressar una variable en funció de l'altra.
\item \textit{Funció a optimitzar.} Escriu la longitud de tanca i multiplica-la per $2{,}5\,\euro/\text{m}$: obtindràs el cost $C$ en funció d'una sola variable.
\item \textit{Optimitza.} Deriva $C$, resol $C'(x) = 0$ i comprova que el valor obtingut és un mínim.
\end{itemize}

\medskip
\textbf{Atenció:} només hi ha \emph{tres} costats de tanca: el quart el fa el mur del jardí.

\medskip
La resposta final ha d'incloure les \emph{dues} dimensions del tancat \emph{i} el cost mínim.}

\subsection*{\textit{b)} \normalfont Si manteniu la forma rectangular o quadrada del tancat i feu que un dels vèrtexs del jardí coincideixi amb un vèrtex del tancat, quants euros us podeu estalviar? Raoneu com posaríeu el tancat i justifiqueu amb càlculs matemàtics les dimensions de la vostra proposta. \hfill\textnormal{[0,75~punts]}}

\pista{Si poses un vèrtex del tancat sobre un vèrtex del jardí, aprofites \emph{dos} murs (els dos costats que hi conflueixen). Només et cal construir \emph{dos} costats de tanca.
\begin{itemize}
\item La condició de l'àrea és la mateixa: $x \cdot y = 8$.
\item Quina és ara la longitud de la tanca? Escriu el nou cost.
\item Optimitza igual que abans per obtenir les dimensions del tancat.
\item Compara el cost mínim amb el de l'apartat \textit{a)}: la diferència és l'estalvi en euros.
\end{itemize}}

\end{document}
